% year: תשפ"ה
% semester: א
% moed: לדוגמה
% number: 2א
% points: 11
% categories: func-analysis, func-limits
% source: exams/מבחנים/תשפו/exam_example 2025.pdf

\begin{question}
(11 נק') מצאו את כל האסימפטוטות של הפונקציה
\[ f(x) = \frac{2x^2 - 3x + lan(x)}{x} \]
נמקו.
\end{question}

\begin{hint}
$lan(x)$ הוא כנראה $\ln(x)$. קבעו את תחום ההגדרה, ובדקו אסימפטוטה אנכית ב-$x = 0^+$ ואסימפטוטה משופעת $y = ax + b$ כאשר $x \to +\infty$.
\end{hint}

\begin{solution}
\textbf{הערה:} בטופס הבחינה מופיע $lan(x)$, וזו כנראה טעות דפוס של $\ln(x)$. נפתור עבור
\[ f(x) = \frac{2x^2 - 3x + \ln x}{x} = 2x - 3 + \frac{\ln x}{x}. \]
\textbf{תחום הגדרה:} $x > 0$ (בגלל $\ln x$; ממילא $x \ne 0$). $f$ רציפה ב-$(0, \infty)$.

\textbf{אסימפטוטה אנכית.} הנקודה היחידה החשודה היא קצה התחום $x = 0$. כאשר $x \to 0^+$: $\ln x \to -\infty$ ו-$\frac1x \to +\infty$, ולכן $\frac{\ln x}{x} = \ln x\cdot\frac1x \to -\infty$, ואילו $2x - 3 \to -3$. לכן
\[ \lim_{x\to0^+} f(x) = -\infty, \]
ו-$x = 0$ אסימפטוטה אנכית (מימין). בכל נקודה אחרת $f$ רציפה ולכן אין שם אסימפטוטה אנכית.

\textbf{אסימפטוטה משופעת.} התחום אינסופי רק ב-$+\infty$. נחשב
\[ a = \lim_{x\to\infty}\frac{f(x)}{x} = \lim_{x\to\infty}\left( 2 - \frac3x + \frac{\ln x}{x^2} \right) = 2, \]
\[ b = \lim_{x\to\infty}\bigl(f(x) - 2x\bigr) = \lim_{x\to\infty}\left( -3 + \frac{\ln x}{x} \right) = -3, \]
כאשר השתמשנו ב-$\lim_{x\to\infty}\frac{\ln x}{x} = 0$ (לפי לופיטל: $\lim \frac{1/x}{1} = 0$), ומכאן גם $\frac{\ln x}{x^2} \to 0$. לכן $y = 2x - 3$ אסימפטוטה משופעת ב-$+\infty$, ואין אסימפטוטה אופקית.

\textbf{תשובה:} אסימפטוטה אנכית $x = 0$, ואסימפטוטה משופעת $y = 2x - 3$ כאשר $x \to +\infty$.
\end{solution}
